% Generated draft: requires mathematical and source review.
\documentclass[11pt]{article}
\usepackage[margin=25mm]{geometry}
\usepackage{fontspec}
\setmainfont{Latin Modern Roman}
\usepackage{amsmath,amssymb,mathtools,unicode-math}
\setmathfont{Latin Modern Math}
\usepackage{enumitem,needspace,xurl,hyperref}
\setlength{\parindent}{0pt}
\setlength{\parskip}{5pt}
\setlength{\emergencystretch}{3em}
\newcommand{\N}{\mathbb{N}}
\newcommand{\Z}{\mathbb{Z}}
\newcommand{\Q}{\mathbb{Q}}
\newcommand{\R}{\mathbb{R}}
\newcommand{\C}{\mathbb{C}}
\newcommand{\F}{\mathbb{F}}
\newcommand{\A}{\mathbb{A}}
\newcommand{\PP}{\mathbb{P}}
\newcommand{\E}{\mathbb{E}}
\newcommand{\eps}{\varepsilon}
\newcommand{\dd}{\,\mathrm{d}}
\begin{document}
\section*{UnsolvedMath ID: 20001888}
\textbf{Source problem code:} catalogue entry
\subsection{Definitions and mathematical statement}
The source's Conjecture 5.2.1 says that compact negatively curved Riemannian manifolds with \(C^0\)-conjugate geodesic flows must be isometric, and then breaks off at ``As with''. Following the aid, interpret compact as closed and connected and conjugacy as time preserving. Thus, for smooth metrics on \(M,N\), a homeomorphism
\[
 H:SM\longrightarrow SN,\qquad
 H\circ\phi_g^t=\phi_h^t\circ H\quad(t\in\mathbb R)
\]
is assumed. This is stronger than an orbit equivalence allowing time change. The conjectural conclusion is an isometry between \((M,g)\) and \((N,h)\).

For a free homotopy class \([a]\), define the marked length spectrum by
\[
 \ell_g([a])=\inf\{L_g(\gamma):\gamma\text{ is a closed loop in }[a]\}.
\]
The paired aid reports the following reduction without assuming that \(H\) covers a base map. If the common dimension is at least three, \(H\) induces an outer isomorphism
\(\theta:\pi_1(M)\to\pi_1(N)\) satisfying
\[
 \ell_g([a])=\ell_h([\theta(a)])
 \qquad\text{for every conjugacy class }[a].
\]
The unit-sphere fibers are simply connected in this range, so the unit-tangent-bundle fundamental groups recover the base groups. The report explicitly treats primitive periodic orbits and their iterates; preserving only an unmarked list of periods would be weaker. It reports the same marked conclusion for oriented negatively curved surfaces, using the characteristic central fiber subgroup of the unit tangent bundle. It does not supply this surface argument for nonorientable surfaces.

A separate lemma in the aid assumes that a phase-space conjugacy covers a continuous map \(F:M\to N\):
\[
 \pi_N\circ H=F\circ\pi_M.
\]
It reports, without a curvature assumption, that \(F\) is a surjective Riemannian covering, hence an isometry if it is injective or a homotopy equivalence. This extra base-map hypothesis must not be inferred merely from the previously recovered fundamental-group marking.

The first result reduces the original conjecture to the relevant marked-length-spectrum rigidity problem; it does not prove that problem in all dimensions. The conditional base-covering lemma settles its stated additional setting and is not a claimed general descent of every phase-space conjugacy.
\subsection{Short English statement}
The conjecture asks whether a time-preserving topological conjugacy of negatively curved geodesic flows forces an isometry. The aid recovers a marking and equal marked length spectra in dimensions at least three and for oriented surfaces. It also gives an isometry criterion when the conjugacy already covers a base map; neither reduction is asserted to settle the unrestricted conjecture.
\subsection{Sources}
\begin{itemize}
\item Local source record: catalogue entry, catalogue entry; canonical \texttt{aim-geometry-notes.json}, \texttt{notes[225]}, source item 5.2, extracted Conjecture 5.2.1. Title: ``Recovering the marking from a flow conjugacy.''
\item AIM workshop ``Geodesics,'' supplied reference \url{https://aimath.org/pastworkshops/geodesicsproblems.pdf}. The original text is truncated after ``As with''. Original remarks and literature fields are empty; this draft uses the supplied text without accessing a PDF.
\item Supplied research attempt 1: marking recovery through sphere fibers, oriented-surface central-fiber argument and conditional Riemannian-covering lemma. No authoritative external bibliographic citation is supplied.
\item Registry dated 14 August 2026: ``partially\_solved,'' ``PARTIAL-PROGRESS,'' statement status ``reconstructed\_unverified''; attempt result ``reduction,'' research-run status open, rigor ``proved,'' novelty confidence low. The machine-generated, non-peer-reviewed aid requires independent verification before citation; the verification field only introduces the canonical record.
\end{itemize}
\end{document}
